\section{Effective systems and the field of section ratios}
\label{sec:completion}

The moduli denominator now allows effective birationality on the base
of a log Calabi--Yau fibration.  We first make two section comparisons
explicit.  They will retain the actual subfield of the function field
through both the fibration and the passage to a good minimal model.

\begin{lemma}[Rounded section comparisons]\label{lem:rounded-comparison}
All varieties in this statement are normal and projective over a field
of characteristic zero.
\begin{enumerate}
\item
Let $f:X\to Z$ be a contraction and let $D$ and $D_Z$ be
$\Q$-Cartier divisors such that
\[
 D+\frac1a\Div(\psi)=f^*D_Z
\]
for an integer $a>0$ and $\psi\in k(X)^*$.  For every positive $\ell$ divisible
by $a$, the spaces of divisorial sections, regarded as rational
functions, satisfy
\begin{equation}\label{eq:section-pullback}
 H^0\bigl(X,\OO_X(\floor{\ell D})\bigr)
 =\psi^{\ell/a}f^*H^0\bigl(Z,\OO_Z(\floor{\ell D_Z})\bigr).
\end{equation}
\item
Let $X\dashrightarrow X'$ be birational, let $D$ and $D'$ be
$\Q$-Cartier divisors, and suppose that on a common resolution
$p:T\to X$, $q:T\to X'$ one has
\[
 p^*D=q^*D'+E,
\]
where $E\geq0$ is $q$-exceptional.  Under the identification of the
function fields, for every integer $\ell>0$ one has
\begin{equation}\label{eq:section-birational}
 H^0\bigl(X,\OO_X(\floor{\ell D})\bigr)
 =H^0\bigl(X',\OO_{X'}(\floor{\ell D'})\bigr).
\end{equation}
\end{enumerate}
\end{lemma}

\begin{proof}
For a rational function $u$, the inequality
$\Div(u)+\floor{\ell D}\geq0$ is equivalent to
$\Div(u)+\ell D\geq0$, since function orders are integral.
Also, effectivity of a $\Q$-Cartier divisor is preserved by surjective
pullback.  It is detected by such a pullback: over the generic point of
each prime divisor downstairs, a local defining parameter pulls back
with positive order along some divisor upstairs.

For (i), divide a section upstairs by $\psi^{\ell/a}$.  The result
$v$ satisfies
\[
 \Div(v)+\ell f^*D_Z\geq0.
\]
It has no pole on the normal projective generic fibre, so is a base
function.  For $v\in k(Z)$ the displayed divisor equals
$f^*(\Div(v)+\ell D_Z)$.  The effectivity comparison just noted proves
both inclusions in~\eqref{eq:section-pullback}.

For (ii), pull an effective divisor $\Div(u)+\ell D$ to $T$ and push
it to $X'$.  Since $q_*E=0$, this gives
$\Div(u)+\ell D'\geq0$.  Conversely, pull the latter inequality to
$T$, add $\ell E$, and push to $X$.  This yields the original
inequality.  The argument uses $\Q$-Cartier pullbacks, so it does not
require $\ell D$ or $\ell D'$ to be Cartier.
\end{proof}

\begin{proposition}[Effective degree on the base]
\label{prop:effective-fibration}
Fix $d\geq1$ and a finite set $\Phi\subset[0,1]\cap\Q$.
There is an integer $N=N(d,\Phi)>0$ with the following property.
Let $f:X\to Z$ be a contraction of normal projective complex varieties
with $\dim X\leq d$ and $\dim Z>0$, and let $(X,B)$ be lc with
$B\geq0$, coefficients in $\Phi$, and $D=K_X+B$ $\Q$-Cartier.
Suppose that $D\simq f^*L$ for a big $\Q$-Cartier divisor $L$ on $Z$.
Then $|\floor{ND}|$ is nonempty, and its section ratios generate
exactly $f^*\C(Z)$ inside $\C(X)$.  This is also the field $K(D)$
generated by section ratios in all positive degrees.
\end{proposition}

\begin{proof}
Use Proposition~\ref{prop:exact-presentation} to obtain
\eqref{eq:exact-presentation}, and choose a smooth projective
determination $\tau:W\to Z$ of its moduli b-divisor.  After a further
resolution we may suppose that the strict transform of $B_Z$ and the
exceptional divisor of $\tau$ have SNC support.  Theorem~
\ref{thm:moduli-denominator} gives a uniform $p$ for which $pM_W$ is
nef Cartier.  Let
\[
 A=\tau_*^{-1}B_Z+\operatorname{Exc}(\tau)_{\mathrm{red}}.
\]
Then $(W,A)$ is log smooth and lc.  Its coefficients lie in the fixed
DCC set for the discriminant, enlarged by $1$.  Moreover,
\begin{equation}\label{eq:base-effective-exceptional}
 K_W+A+M_W=\tau^*D_Z+E,\qquad E=A-B_W\geq0,
\end{equation}
where $E$ is exceptional: at nonexceptional primes $A$ and $B_W$
agree, and at exceptional primes $A$ has coefficient one whereas
$B_W$ has coefficient at most one.

The divisor in~\eqref{eq:base-effective-exceptional} is big.
The effective birationality theorem for polarized pairs
\cite[Theorem~1.3]{BZ} now applies to $(W,A)$ and $M_W$: the pair is
projective lc, its boundary coefficients belong to a fixed DCC set,
$pM_W$ is nef Cartier, and the adjoint sum is big.  Consequently
\[
 \left|\floor{n(K_W+A+M_W)}\right|
\]
is birational for every $n$ divisible by an integer depending only on
$\dim W$, the DCC set, and $p$.  The notation in that theorem uses
round down, as in the displayed system.

Pushing the rational section inequalities to $Z$ shows that this
birational system is a subsystem of $|\floor{nD_Z}|$ under the
identification $\C(W)=\C(Z)$.  Thus the latter system is nonempty
and its ratios generate $\C(Z)$.  Choose a common such $n=N$ over
$1\leq\dim Z\leq d$, also divisible by $p_0$.
Equation~\eqref{eq:section-pullback} then proves the asserted
nonemptiness and ratio-field equality upstairs.

It remains to compare with all degrees.  If $s,s_0$ are nonzero
sections in degree $\ell$, their ratio can be represented in every
multiple degree $b\ell$ as
\begin{equation}\label{eq:ratio-multiple}
 \frac{s}{s_0}=\frac{s\,s_0^{b-1}}{s_0^b}.
\end{equation}
These products are sections in degree $b\ell$ because
$b\floor{\ell D}\leq\floor{b\ell D}$.
Choose $b$ so that $p_0\mid b\ell$ and apply
\eqref{eq:section-pullback}.  Every such ratio belongs to
$f^*\C(Z)$, proving $K(D)=f^*\C(Z)$.
\end{proof}

The proof of the main theorem first treats the ground field $\C$.
For the final change of ground field we use the following elementary
observation, which keeps track of equality of subfields rather than only
the associated rational maps.

\begin{lemma}[Descent of the field of ratios]\label{lem:field-descent}
Let $k\subset k'$ be an extension of algebraically closed fields, let
$X$ be a normal integral projective variety over $k$, and let $D$ be a
$\Q$-divisor on $X$.  Write $X'=X_{k'}$ and $D'=D_{k'}$.
For each $\ell>0$ having nonzero sections, let $F_\ell\subset k(X)$
be the field generated over $k$ by ratios of sections of
$\OO_X(\floor{\ell D})$, and define $F'_\ell$ similarly on $X'$.
Then
\[
 F'_\ell=k'F_\ell,\qquad K(D')=k'K(D),
\]
where the composita are taken in $k'(X')$.  In particular, for a fixed
$m>0$, nonemptiness of $|\floor{mD}|$ and the equality
$F_m=K(D)$ hold if and only if they hold after extension to $k'$.
\end{lemma}

\begin{proof}
The rounded divisorial sheaves commute with these field extensions.
One may check this on the smooth big open set where the prime divisors
are Cartier, and then use reflexive extension across its complement.
Base change for global sections gives
\begin{equation}\label{eq:section-base-change}
 H^0\bigl(X',\OO_{X'}(\floor{\ell D'})\bigr)
 =k'\otimes_k H^0\bigl(X,\OO_X(\floor{\ell D})\bigr).
\end{equation}
Choose a nonzero section $s_0$ over $k$.  Dividing any section over
$k'$ by $s_0$ expresses it as a finite $k'$-linear combination of
ratios over $k$.  This proves $F'_\ell=k'F_\ell$; taking all degrees
proves the analogous assertion for $K(D)$.

For descent of the equality, we use the following intersection fact.
If $F$ is an intermediate field $k\subset F\subset k(X)$, then
\begin{equation}\label{eq:field-intersection}
 k(X)\cap k'F=F
\end{equation}
inside $k'(X')$.  Indeed $X$ is geometrically integral, so
$k(X)\otimes_k k'$ injects into $k'(X')$.  If $u\in k(X)\cap k'F$,
write
\[
 u=\frac{\sum_{i=1}^a c_i a_i}{\sum_{i=1}^a c_i b_i},
 \qquad a_i,b_i\in F,\quad c_i\in k',
\]
where the $c_i$ are linearly independent over $k$ and the denominator
is nonzero.  Such an expression is obtained by collecting a finite
list of constants into a $k$-basis.  After clearing the denominator,
injectivity of the tensor product map and linear independence give
$ub_i=a_i$ for every $i$.  Some $b_i$ is nonzero, so $u\in F$.
Applying~\eqref{eq:field-intersection} to $F=F_m$ proves descent of
$F'_m=K(D')$.  The forward implication follows from the compositum
identities, and nonemptiness follows in either direction from
\eqref{eq:section-base-change}.
\end{proof}

\begin{proof}[Proof of Theorem~\ref{thm:main}]
First let the ground field be $\C$, and put $D=K_X+B$.
The hypothesis $\kappa(X,D)\geq0$ implies that $D$ is
pseudo-effective.  Take a crepant $\Q$-factorial dlt modification
$\mu:(Y,\Delta)\to(X,B)$.  Its adjoint is $\mu^*D$, and its
boundary coefficients belong to $\Phi\cup\{1\}$.
Theorem~\ref{thm:good-models} and termination with ample scaling
\cite[Theorem~A]{TX} give a terminating $(K_Y+\Delta)$-MMP to a
$\Q$-factorial dlt pair $(X',B')$.  The nef case of the good-model
theorem makes its adjoint semiample.  The MMP extracts no divisors,
so the coefficients of $B'$ still belong to $\Phi\cup\{1\}$.
With compatible canonical divisors, crepancy of $\mu$ and the
discrepancy comparison for this MMP give, on a common resolution,
\[
 p^*(K_X+B)=q^*(K_{X'}+B')+E,
 \qquad E\geq0\quad\text{$q$-exceptional};
\]
see the minimal-model discrepancy comparison and negativity lemma
in~\cite{KM}.  Lemma~\ref{lem:rounded-comparison}(ii) identifies all
rounded section spaces as subspaces of the common function field.
It therefore suffices to prove the result on $(X',B')$.

The adjoint $D'=K_{X'}+B'$ is semiample.  Its semiample contraction
$f:X'\to Z$ satisfies $D'\simq f^*L$ with $L$ an ample
$\Q$-Cartier divisor on the normal projective base.  If $\dim Z>0$,
Proposition~\ref{prop:effective-fibration}, with coefficient set
$\Phi\cup\{1\}$, gives a uniform nonempty degree whose ratios
generate the entire field $K(D')$.
If $Z$ is a point, then $D'\simq0$, and
Theorem~\ref{thm:normal-index} gives a uniform principal multiple.
Its system is nonempty and all ratios in that degree are constants.
For any other degree, pass to a common multiple by
\eqref{eq:ratio-multiple}; ratios in the latter degree are again
constants.  Hence $K(D')=\C$ in this case.  A common multiple of the
two uniform degrees works in both cases, again by
\eqref{eq:ratio-multiple}.  Denote it by $m(d,\Phi)$.  The rounded
comparison transfers the conclusion to $X$.

Now let $k$ be any algebraically closed field of characteristic zero.
Descend $X$, its prime boundary components and coefficients, the
canonical and $\Q$-Cartier data, and a log resolution to a finitely
generated subfield of $k$.  Let $k_0\subset k$ be the algebraic
closure of that field inside $k$.  Enlarging the finitely generated
field if necessary, the descended variety $X_0$ is geometrically
integral and normal, its pair $(X_0,B_0)$ is lc, and its adjoint
$D_0$ is $\Q$-Cartier.  These conditions can be checked after the
faithfully flat extension to $k$, using the chosen resolution for
the discrepancy inequalities.  The field $k_0$ embeds into $\C$.

Equation~\eqref{eq:section-base-change} shows that nonzero sections in
some positive degree, and hence nonnegative Kodaira dimension, descend
from $k$ to $k_0$ and persist after extension to $\C$.
The complex case gives the conclusion in the same integer
$m(d,\Phi)$ for $(X_0,B_0)_\C$.
Lemma~\ref{lem:field-descent} first descends the nonemptiness and exact
equality of ratio fields to $k_0$, and then extends them to $k$.
This proves the theorem over the stated ground field.
\end{proof}
